% GENERATED FILE. Edit canonical lessons, then run npm run generate:books.
% canonical-source-hash: fnv1a64:2bc8d26673f9140d
% canonical-lessons: IT-C316-candelina, IT-C316-festeggiato, IT-C316-incartare, IT-C316-fiocco, IT-C316-palloncino

\chapter{La torta di compleanno}
\label{ch:it-la-torta-di-compleanno}
\hlchaptermodality{316}

\begin{chapteropening}
\textbf{By the end of this chapter:} \emph{I can describe a birthday party: the candles, the birthday person, wrapping presents, bows and balloons.}

Complete the last lesson of chapter 316: I can describe a birthday party: the candles, the birthday person, wrapping presents, bows and balloons.
\end{chapteropening}

\section[la candelina]{la candelina — the birthday candle}

\label{lesson:IT-C316-candelina}



\begin{quote}
Before the new one: say the Italian for best wishes, then the Italian for to wish.
\end{quote}

\subsection*{You'll want to know: la candelina}
\textbf{la candelina} — \textquotedblleft{}the birthday candle\textquotedblright{}.

\begin{quote}
Sulla torta ci sono trenta candeline.
\end{quote}

\begin{quote}
\emph{There are thirty candles on the cake.}
\end{quote}

A little \emph{candela}, ‘candle’.

\begin{cousinweb}
\emph{candelina} is a little \emph{candela} \textquotedblleft{}candle\textquotedblright{}, inherited from Latin \emph{candēla}, from \emph{candēre} \textquotedblleft{}to shine white\textquotedblright{} (Lewis \& Short).
\end{cousinweb}

\begin{grammarlens}[title={What you've built}]
One more word for telling, explaining and describing at B1.
\end{grammarlens}

\subsection*{Your turn}
\begin{itemize}
\raggedright
  \item \emph{Say it:} \emph{la candelina}
  \item \emph{Say it:} \emph{la candelina}, once more
  \item \emph{Say it:} a sentence of your own with \emph{la candelina}
  \item \emph{Recall:} say the Italian for best wishes, then the Italian for to wish, then say \emph{la candelina} again
\end{itemize}

\begin{tcolorbox}[breakable,colback=teal!4,colframe=teal!35!black,title={Before you move on}]
What does la candelina mean? (\textquotedblleft{}The birthday candle\textquotedblright{}.) Say it once more.
\end{tcolorbox}
\section[il festeggiato]{il festeggiato — the person being celebrated, the birthday boy or girl}

\label{lesson:IT-C316-festeggiato}



\begin{quote}
Before the new one: say the Italian for to wish, then the Italian for the birthday candle.
\end{quote}

\subsection*{You'll want to know: il festeggiato}
\textbf{il festeggiato} — \textquotedblleft{}the person being celebrated, the birthday boy or girl\textquotedblright{}.

\begin{quote}
Il festeggiato ha aperto i regali.
\end{quote}

\begin{quote}
\emph{The birthday boy opened his presents.}
\end{quote}

Built on \emph{festeggiare}; a woman is \emph{la festeggiata}.

\begin{cousinweb}
\emph{festeggiato} is the past participle of \emph{festeggiare} \textquotedblleft{}to celebrate\textquotedblright{}, built on \emph{festa}, from Latin \emph{fēsta} \textquotedblleft{}feast days\textquotedblright{} (Lewis \& Short).
\end{cousinweb}

\begin{grammarlens}[title={What you've built}]
One more word for telling, explaining and describing at B1.
\end{grammarlens}

\subsection*{Your turn}
\begin{itemize}
\raggedright
  \item \emph{Say it:} \emph{il festeggiato}
  \item \emph{Say it:} \emph{il festeggiato}, once more
  \item \emph{Say it:} a sentence of your own with \emph{il festeggiato}
  \item \emph{Recall:} say the Italian for to wish, then the Italian for the birthday candle, then say \emph{il festeggiato} again
\end{itemize}

\begin{tcolorbox}[breakable,colback=teal!4,colframe=teal!35!black,title={Before you move on}]
What does il festeggiato mean? (\textquotedblleft{}The person being celebrated, the birthday boy or girl\textquotedblright{}.) Say it once more.
\end{tcolorbox}
\section[incartare]{incartare — to wrap (in paper)}

\label{lesson:IT-C316-incartare}



\begin{quote}
Before the new one: say the Italian for the birthday candle, then the Italian for the person being celebrated.
\end{quote}

\subsection*{You'll want to know: incartare}
\textbf{incartare} — \textquotedblleft{}to wrap (in paper)\textquotedblright{}.

\begin{quote}
Può incartare il regalo, per favore?
\end{quote}

\begin{quote}
\emph{Could you wrap the present, please?}
\end{quote}

Built on \emph{carta}, ‘paper’.

\begin{cousinweb}
\emph{incartare} is built on \emph{carta} \textquotedblleft{}paper\textquotedblright{}, from Latin \emph{charta} (Lewis \& Short).
\end{cousinweb}

\begin{grammarlens}[title={What you've built}]
One more word for telling, explaining and describing at B1.
\end{grammarlens}

\subsection*{Your turn}
\begin{itemize}
\raggedright
  \item \emph{Say it:} \emph{incartare}
  \item \emph{Say it:} \emph{incartare}, once more
  \item \emph{Say it:} a sentence of your own with \emph{incartare}
  \item \emph{Recall:} say the Italian for the birthday candle, then the Italian for the person being celebrated, then say \emph{incartare} again
\end{itemize}

\begin{tcolorbox}[breakable,colback=teal!4,colframe=teal!35!black,title={Before you move on}]
What does incartare mean? (\textquotedblleft{}To wrap (in paper)\textquotedblright{}.) Say it once more.
\end{tcolorbox}
\section[il fiocco]{il fiocco — the bow (ribbon); the flake}

\label{lesson:IT-C316-fiocco}



\begin{quote}
Before the new one: say the Italian for the person being celebrated, then the Italian for to wrap.
\end{quote}

\subsection*{You'll want to know: il fiocco}
\textbf{il fiocco} — \textquotedblleft{}the bow (ribbon); the flake\textquotedblright{}.

\begin{quote}
Sul pacco c'è un fiocco rosso.
\end{quote}

\begin{quote}
\emph{There is a red bow on the parcel.}
\end{quote}

\emph{Fiocchi di neve} — snowflakes.

\begin{cousinweb}
\emph{fiocco} is inherited from Latin \emph{floccus} \textquotedblleft{}a tuft of wool\textquotedblright{} (EtymDB).
\end{cousinweb}

\begin{grammarlens}[title={What you've built}]
One more word for telling, explaining and describing at B1.
\end{grammarlens}

\subsection*{Your turn}
\begin{itemize}
\raggedright
  \item \emph{Say it:} \emph{il fiocco}
  \item \emph{Say it:} \emph{il fiocco}, once more
  \item \emph{Say it:} a sentence of your own with \emph{il fiocco}
  \item \emph{Recall:} say the Italian for the person being celebrated, then the Italian for to wrap, then say \emph{il fiocco} again
\end{itemize}

\begin{tcolorbox}[breakable,colback=teal!4,colframe=teal!35!black,title={Before you move on}]
What does il fiocco mean? (\textquotedblleft{}The bow (ribbon); the flake\textquotedblright{}.) Say it once more.
\end{tcolorbox}
\section[il palloncino]{il palloncino — the balloon}

\label{lesson:IT-C316-palloncino}



\begin{quote}
Before the new one: say the Italian for to wrap, then the Italian for the bow.
\end{quote}

\subsection*{You'll want to know: il palloncino}
\textbf{il palloncino} — \textquotedblleft{}the balloon\textquotedblright{}.

\begin{quote}
Abbiamo comprato cinquanta palloncini.
\end{quote}

\begin{quote}
\emph{We bought fifty balloons.}
\end{quote}

A little \emph{pallone}, ‘ball’.

\begin{cousinweb}
\emph{palloncino} is a little \emph{pallone} \textquotedblleft{}big ball\textquotedblright{}, built on \emph{palla}, which EtymDB traces to a Longobard word for a ball.
\end{cousinweb}

\begin{grammarlens}[title={What you've built}]
One more word for telling, explaining and describing at B1.
\end{grammarlens}

\subsection*{Your turn}
\begin{itemize}
\raggedright
  \item \emph{Say it:} \emph{il palloncino}
  \item \emph{Say it:} \emph{il palloncino}, once more
  \item \emph{Say it:} a sentence of your own with \emph{il palloncino}
  \item \emph{Recall:} say the Italian for to wrap, then the Italian for the bow, then say \emph{il palloncino} again
\end{itemize}

\begin{tcolorbox}[breakable,colback=teal!4,colframe=teal!35!black,title={Before you move on}]
What does il palloncino mean? (\textquotedblleft{}The balloon\textquotedblright{}.) Say it once more.
\end{tcolorbox}
